\subsection{管理维护模块}

\paragraph*{业务流程}

\paragraph*{用例图}

\subsubsection{功能点——实体添加（FR-GLWH-0001）}

\paragraph*{简述}

管理员可通过管理端界面直接添加用户、图像等实体，新增实体可在用户端显示或使用。

\paragraph*{前提条件}

管理员已登录系统，且待添加实体的数据符合系统规范（如格式、字段完整性）。

\paragraph*{主要流程}

\begin{enumerate}
    \item 管理员进入“实体管理”页面，选择“添加实体”。
    \item 填写实体信息（如用户账号、图像元数据等），上传相关文件。
    \item 系统校验数据合法性后保存至数据库。
\end{enumerate}

\paragraph*{后继结果}

新增实体成功显示在用户端对应模块。若因各种原因导致失败，系统提示错误原因并要求重试或者重新输入。

\subsubsection{功能点——实体修改（FR-GLWH-0002）}

\paragraph*{简述}

管理员可修改用户、图像等实体的信息，修改后同步至用户端。

\paragraph*{前提条件}

目标实体已存在于数据库中，且修改后的数据符合系统规范。

\paragraph*{主要流程}

\begin{enumerate}
    \item 管理员进入“实体管理”页面，找到目标实体
    \item 选择“编辑”操作，修改相关信息（如用户权限、图像分类标签）。
    \item 提交修改后，系统更新数据库并同步至用户端。
\end{enumerate}

\paragraph*{后继结果}

用户端显示修改后的实体信息。若因各种与原因导致失败，系统提示错误原因并要求重试或者重新修改。

\subsubsection{功能点——实体删除（FR-GLWH-0003）}

\paragraph*{简述}

管理员可删除指定的用户、图像等实体，删除后实体不可恢复且不会在用户端显示。

\paragraph*{前提条件}

目标实体存在于数据库中且删除实体不会引发错误。

\paragraph*{主要流程}

\begin{enumerate}
    \item 管理员进入“实体管理”页面，搜索目标实体。
    \item 选择“删除”操作，系统提示确认对话框。
    \item 确认后，系统将实体标记为“已删除”，并移出用户端展示列表。
\end{enumerate}

\paragraph*{后继结果}

实体在用户端不可见且无法被检索。若删除失败，则反馈失败原因并要求重试或放弃。

\subsubsection{功能点——实体检索（FR-GLWH-0004）}

\paragraph*{简述}

管理员可通过模糊检索快速定位用户、图像等实体。

\paragraph*{前提条件}

管理员登录系统。

\paragraph*{主要流程}

\begin{enumerate}
    \item 管理员进入“实体管理”页面，在搜索栏输入关键词（如用户昵称、图像ID）；当对图像进行检索时，还可按照图像标签分类进行检索。
    \item 系统实时匹配数据库中的实体名称，展示结果列表。
    \item 管理员可点击条目查看详情或进行进一步操作。
\end{enumerate}

\paragraph*{后继结果}

若有匹配结果，则以条目形式呈现；若无匹配结果，系统提示“未找到相关实体”。

\subsubsection{功能点——实体分类（FR-GLWH-0005）}

\paragraph*{简述}

管理员可按学科、造假程度等标签对图像进行分类管理。

\paragraph*{前提条件}

管理员登陆系统。

\paragraph*{主要流程}

\begin{enumerate}
    \item 管理员进入“实体管理”页面，选择图像分类。
    \item 系统提供分类情况。
\end{enumerate}

\paragraph*{后继结果}

图像按照分类标准进行有序展示，便于后续统计和查询。

\subsubsection{功能点——数据信息统计（FR-GLWH-0006）}

\paragraph*{简述}

管理员可查看用户、图像、检测任务等维度的统计报表。

\paragraph*{前提条件}

系统已积累一定量的业务数据。

\paragraph*{主要流程}

\begin{enumerate}
    \item 管理员进入“数据统计”页面，选择统计维度（如“用户活跃度”）。
    \item 系统生成可视化图表（如柱状图、饼图）及数据表格并直接显示。
    \item 管理员可导出报表为Excel或PDF格式。
\end{enumerate}

\paragraph*{后继结果}

管理员获得直观数据统计。可导出文件保存至本地，供后续分析使用。

\subsubsection{功能点——操作记录统计（FR-GLWH-0007）}

\paragraph*{简述}

管理员可查看系统日志，包括用户操作、检测任务记录等。

\paragraph*{前提条件}

管理员登录系统。

\paragraph*{主要流程}

\begin{enumerate}
    \item 管理员进入“系统日志”页面，选择日志类型（如“用户登录日志”）。
    \item 系统按时间倒序展示日志条目，支持按关键词筛选。
    \item 管理员可导出日志或标记异常事件。
\end{enumerate}

\paragraph*{后继结果}

管理员完成系统审计或故障排查，对异常日志进行处理。

\subsubsection{功能点——超级管理员（FR-GLWH-0008）}

\paragraph*{简述}

超级管理员可管理普通管理员账号，包括添加、删除及权限分配。

\paragraph*{前提条件}

当前用户为超级管理员角色，且目标账号目前可以更改权限。

\paragraph*{主要流程}

\begin{enumerate}
    \item 超级管理员进入“管理员管理”页面。
    \item 选择“添加管理员”，填写账号信息并分配权限。
    \item 或选择“删除管理员”，确认后移除目标账号权限。
\end{enumerate}

\paragraph*{后继结果}

新增管理员账号可正常登录并行使权限，而被删除的管理员账号立即失效，无法操作系统。

\subsubsection{功能点——用户举报审核（FR-GLWH-0009）}

\paragraph*{简述}

管理员处理用户举报信息，并根据结果执行警告、封号等操作。

\paragraph*{前提条件}

用户已提交举报并附证据，该举报内容还未被处理。

\paragraph*{主要流程}

\begin{enumerate}
    \item 管理员进入“举报处理”页面，查看举报详情及证据。
    \item 核实后选择处理方式（如“警告”“封号”），填写处理说明。
    \item 系统向举报方和被举报方发送处理结果通知。
\end{enumerate}

\paragraph*{后继结果}

举报方收到处理反馈，被举报方账号状态更新（如受限）。处理记录保存至日志，供后续复查。

\subsubsection{功能点——学术资源管理（FR-GLWH-0010）}

\subsubsection{功能点——论文检测日志管理（FR-GLWH-0011）}

\subsubsection{功能点——Review检测日志管理（FR-GLWH-0012）}

\subsubsection{功能点——模型管理与配置（FR-GLWH-0013）}